\begin{proposition}[The integral quintic isogeny]
\label{prop:quintic-five-isogeny}
For the maps in \eqref{eq:maps}, put
\[
 \Phi=(f_{Z,*},f_{X,*},f_{Y,*}):\Jac(X_5)\longrightarrow A^3,
 \qquad
 \Psi=f_Z^*+f_X^*+f_Y^*:A^3\longrightarrow\Jac(X_5).
\]
Then
\[
 \Phi\Psi=[5]_{A^3},\qquad \Psi\Phi=[5]_{\Jac(X_5)}.
\]
In particular the kernels are killed by five, and
$T_3\Jac(X_5)\simeq(T_3A)^{\oplus3}$ as integral $G_\Q$-modules.
\end{proposition}
\begin{proof}
Each $f_i$ is the quotient by an order-five diagonal subgroup $H_i$.
For example the first map is invariant under
$(x,y)\mapsto(\zeta_5x,\zeta_5^{-1}y)$; its function field has
degree five because $x^5=v$ and $y=-u/x$.
The other two cases follow by cyclic permutation.
Consequently $f_{i,*}f_i^*=[5]$ and
$f_i^*f_{i,*}=\sum_{h\in H_i}h$.

The characters of the holomorphic differentials
$\omega_{a,b}=x^{a-1}y^{b-5}dx$ have $a,b\ge1$ and $a+b\le4$.
The three invariant pairs, in the order $Z,X,Y$, are
\[
 \{(1,1),(2,2)\},\qquad
 \{(1,2),(3,1)\},\qquad
 \{(1,3),(2,1)\}.
\]
They are disjoint and exhaust the six characters, as also follows
from \eqref{eq:differentials}. The $H_i$-trace is therefore five
on its own invariant subspace and zero on the other two. This
proves that $f_{i,*}f_j^*=0$ for $i\ne j$ and that the sum of
the three trace operators is $[5]$. These are equalities of
homomorphisms: a homomorphism of abelian varieties in characteristic
zero with zero differential is zero. More explicitly, the product
of distinct trace operators is zero; $f_{j,*}$ is surjective and
$f_i^*$ has finite kernel, so the mixed composition itself is zero.
The two displayed compositions
follow. Since five is a unit in $\Z_3$, the integral inverse of
$T_3\Phi$ is $5^{-1}T_3\Psi$.
\end{proof}

\begin{lemma}[An integral splitting of the quintic boundary]
\label{lem:quintic-boundary-split}
Let $D_5=X_5\setminus U_5$ and
$B_5=\Z_3(1)[D_5(\bar\Q)]/\Z_3(1)\sum_{Q\in D_5}e_Q$.
The sequence
\[
 0\longrightarrow B_5\longrightarrow
 \pi_1(U_{5,\bar\Q})^{\ab,(3)}
 \longrightarrow T_3\Jac(X_5)\longrightarrow0
\]
splits as a sequence of integral $G_\Q$-modules.
\end{lemma}
\begin{proof}
The generalized Jacobian extension by the boundary torus is
classified by the map from degree-zero divisors supported on $D_5$
to $\Jac(X_5)$, using the principal polarization.
For coordinate cusps $P,Q$, the quotient of their tangent-line
equations has divisor $5(Q-P)$. Thus this classifying map is killed
by five. Pullback of the extension by $[5]$ is trivial. Its splitting
is unique because a homomorphism from an abelian variety to a torus
is zero, and hence the splitting descends to $\Q$.
The resulting lift of $[5]$, divided by five on the $3$-adic Tate
module, gives the asserted splitting.
\end{proof}

These integral statements strengthen Proposition~\ref{prop:isogeny}.
For the pro-$3$ image argument below, a rational isogeny already
suffices: having pro-$3$ image is a topological property of the
rational Tate representation and is independent of a stable lattice.
